% Please do not change %
\documentclass[11pt]{article} % font size
\usepackage[margin=1in]{geometry} % margins
\usepackage{amsmath,amsthm,amssymb} % for the  common "math symbols"
\usepackage{algorithm}
\usepackage{algpseudocode}
\usepackage[shortlabels]{enumitem} % for enumeration
\usepackage{booktabs} % if you need tables
\usepackage{listings}
\usepackage{color}
\usepackage{quiver}
\usepackage{ulem}
\DeclareMathOperator*{\argmax}{argmax}
\DeclareMathOperator*{\argmin}{argmin}

% \theoremstyle{definition}
\newtheorem{theorem}{Theorem}[section]
\newtheorem{corollary}[theorem]{Corollary}
\newtheorem{lemma}[theorem]{Lemma}
\newtheorem{proposition}[theorem]{Proposition}
\newtheorem{claim}[theorem]{Claim}
\newtheorem{conjecture}[theorem]{Conjecture}
\theoremstyle{definition}
\newtheorem{definition}[theorem]{Definition}
\theoremstyle{definition}
\newtheorem{fact}[theorem]{Fact}
\theoremstyle{definition}
\newtheorem{example}[theorem]{Example}
\newtheorem{note}[theorem]{Note}

\lstset{
  frame=none,
  xleftmargin=2pt,
  stepnumber=1,
  numbers=left,
  numbersep=5pt,
  numberstyle=\ttfamily\tiny\color[gray]{0.3},
  belowcaptionskip=\bigskipamount,
  captionpos=b,
  escapeinside={*'}{'*},
  language=haskell,
  tabsize=2,
  emphstyle={\bf},
  commentstyle=\it,
  stringstyle=\mdseries\rmfamily,
  showspaces=false,
  keywordstyle=\bfseries\rmfamily,
  columns=flexible,
  basicstyle=\small\sffamily,
  showstringspaces=false,
  morecomment=[l]\%,
}
% Feel free to include additional packages (if needed), but make sure it does not override the margin/font requirements.
%
\setlength{\parindent}{0pt} % no indent for new paragraphs.
\setlength{\parskip}{5pt} % distance between paragraphs, which will be used when you have a blank line in your LaTeX code.


%%%%%%%%%%%%%%%%%%%%%%%%%%%%%%%%%%%%%%%%%%%%%%%%%%%%%%%%%%%%%%%%%%%%%%%%%%%%%%
\title{Categories, Proofs, and Processes: \\Lecture 13 Yoneda's Lemma}
%
\author{Wira Azmoon Ahmad}
\date{8 Feb 2024}
%
\begin{document}
%
\maketitle

%You may use the usual \LaTeX\ environments to structure your answers:
%\begin{align}
%        \label{basegnn}
%        \mathbf{h}_u^{(t+1)} = \sigma \bigg( \mathbf{W}_\text{self}^{(t)} \mathbf{h}_u^{(t)}  + \mathbf{W}_\text{neigh}^{(t)} \sum_{v \in N(u)} \mathbf{h}_v^{(t)} + \mathbf{b}^{(t)} \bigg): 
%\end{align}


%\begin{table}[h]
%\caption{A simple table.}
%\centering
%\begin{tabular}[t]{lrrrr}
%\toprule
%Graphs & Can & Be  & Fun & Too \\ 
%\midrule 
%$G_1$ & $1$  & $2$ & $3$ & $4$ \\ 
%$G_2$ & $-1$  & $-2$ & $-3$ & $-4$ \\ 
%$G_3$ & $0.1$  & $0.2$ & $0.3$ & $0.4$ \\ 
%\bottomrule
%\end{tabular}
%\label{tab}
%\end{table}

% A functor $F:\mathcal{C} \rightarrow \mathcal{D}$ is
% \begin{itemize}
%     \item \textbf{faithful} if each map $F_{A,B}:\mathcal{C}(A,B) \rightarrow \mathcal{D}(F\ A, F\ B)$ is injective;
%     \item \textbf{full} if each map $F_{A,B}:\mathcal{C}(A,B) \rightarrow \mathcal{D}(F\ A, F\ B)$ is surjective
% \end{itemize}


\addtocounter{section}{4}
\section{Yoneda's Lemma}
\subsection{Motivation}
There are lots of Hom-functors, one for every $A\in\text{Ob}(\mathcal{C})$.

... but they are all related by natural transformations.
\[t:\mathcal{C}(-,A)\rightarrow\mathcal{C}(-,B) := \{t_C : \mathcal{C}(C,A)\rightarrow\mathcal{C}(C,B)\}\]
For $f:C\rightarrow D$,
% https://q.uiver.app/#q=WzAsNCxbMCwwLCJcXG1hdGhjYWx7Q30oRCxBKSJdLFsxLDAsIlxcbWF0aGNhbHtDfShDLEEpIl0sWzAsMSwiXFxtYXRoY2Fse0N9KEQsQikiXSxbMSwxLCJcXG1hdGhjYWx7Q30oRCxCKSJdLFswLDEsIlxcbWF0aGNhbHtDfShmLEEpIl0sWzAsMiwidF9EIiwyXSxbMiwzLCJcXG1hdGhjYWx7Q30oZixCKSIsMl0sWzEsMywidF9DIl1d
\[\begin{tikzcd}
	{\mathcal{C}(D,A)} & {\mathcal{C}(C,A)} \\
	{\mathcal{C}(D,B)} & {\mathcal{C}(D,B)}
	\arrow["{\mathcal{C}(f,A)}", from=1-1, to=1-2]
	\arrow["{t_D}"', from=1-1, to=2-1]
	\arrow["{\mathcal{C}(f,B)}"', from=2-1, to=2-2]
	\arrow["{t_C}", from=1-2, to=2-2]
\end{tikzcd}\]
The contravariant Hom-functor fixes the first element, the natural transfomration fixes the second element.
\begin{claim}[Ex. 55]
    For $g:A\rightarrow B$,
    \[t:= \mathcal{C}(-,g) := \{t_C = \mathcal{C}(C,g) : \mathcal{C}(C,A)\rightarrow\mathcal{C}(C,B)\}\]
    $\mathcal{C}(-,g)$ is a natural transformation.
\end{claim}
\begin{proof}
    % https://q.uiver.app/#q=WzAsOCxbMCwwLCJcXG1hdGhjYWx7Q30oRCxBKSJdLFszLDAsIlxcbWF0aGNhbHtDfShDLEEpIl0sWzAsMywiXFxtYXRoY2Fse0N9KEQsQikiXSxbMywzLCJcXG1hdGhjYWx7Q30oRCxCKSJdLFsxLDEsImgiXSxbMiwxLCJoXFxjaXJjIGYiXSxbMSwyLCJnXFxjaXJjIGgiXSxbMiwyLCJnXFxjaXJjIGggXFxjaXJjIGYiXSxbMCwxLCJcXG1hdGhjYWx7Q30oZixBKSJdLFswLDIsInRfRCIsMl0sWzIsMywiXFxtYXRoY2Fse0N9KGYsQikiLDJdLFsxLDMsInRfQyJdLFs0LDUsIiIsMCx7InN0eWxlIjp7InRhaWwiOnsibmFtZSI6Im1hcHMgdG8ifX19XSxbNCw2LCIiLDIseyJzdHlsZSI6eyJ0YWlsIjp7Im5hbWUiOiJtYXBzIHRvIn19fV0sWzUsNywiIiwwLHsic3R5bGUiOnsidGFpbCI6eyJuYW1lIjoibWFwcyB0byJ9fX1dLFs2LDcsIiIsMix7InN0eWxlIjp7InRhaWwiOnsibmFtZSI6Im1hcHMgdG8ifX19XV0=
\[\begin{tikzcd}
	{\mathcal{C}(D,A)} &&& {\mathcal{C}(C,A)} \\
	& h & {h\circ f} \\
	& {g\circ h} & {g\circ h \circ f} \\
	{\mathcal{C}(D,B)} &&& {\mathcal{C}(D,B)}
	\arrow["{\mathcal{C}(f,A)}", from=1-1, to=1-4]
	\arrow["{t_D}"', from=1-1, to=4-1]
	\arrow["{\mathcal{C}(f,B)}"', from=4-1, to=4-4]
	\arrow["{t_C}", from=1-4, to=4-4]
	\arrow[maps to, from=2-2, to=2-3]
	\arrow[maps to, from=2-2, to=3-2]
	\arrow[maps to, from=2-3, to=3-3]
	\arrow[maps to, from=3-2, to=3-3]
\end{tikzcd}\]
\end{proof}

\begin{lemma}[Faux-neda]
    There exists a bijection
    \[\phi : \text{Nat}(\mathcal{C}(-,A), \mathcal{C}(-,B))\cong \mathcal{C}(A,B)\]
    \label{faux}
\end{lemma}

\subsection{Yoneda's Lemma}
\begin{lemma}[Yoneda's Lemma]
    For any functor $H:\mathcal{C}^{op} \rightarrow\textbf{Set}$, and $A\in \text{Ob}(\mathcal{C})$,
    there exists a bijection (imagine $H=\mathcal{C}(-,B)$)
    \[\phi : \text{Nat}(\mathcal{C}(-,A), H)\cong H\ A\]
    \label{yo}
\end{lemma}

\begin{proof}
Take a natural transformation $t:\mathcal{C}(-,A)\rightarrow H$. We need an element $u\in H\ A$.
\[t_A:\mathcal{C}(A,A)\rightarrow H\ A\]
Let
\[u:= t_A(1_A)\in H\ A, \quad \phi(t) = t_A(1_A)\]

For the inverse, take some $u\in H\ A$, and build a natural transformation $t$ from it.
\[t_C: \mathcal{C}(C,A) \rightarrow H\ C\]
\[t_C(g) := H\ g\ u\]
Need to show that $t$ is natural. $\forall f:C\rightarrow D$,
% https://q.uiver.app/#q=WzAsOCxbMCwwLCJcXG1hdGhjYWx7Q30oRCxBKSJdLFszLDAsIlxcbWF0aGNhbHtDfShDLEEpIl0sWzAsMywiSFxcIEQiXSxbMywzLCJIXFwgQyJdLFsxLDEsImciXSxbMiwxLCJnXFxjaXJjIGYiXSxbMSwyLCJIXFwgZ1xcIHUiXSxbMiwyLCJIZihIXFwgZ1xcIHUpIl0sWzAsMSwiXFxtYXRoY2Fse0N9KGYsQSkiXSxbMCwyLCJ0X0QiLDJdLFsyLDMsIkhcXCBmIiwyXSxbMSwzLCJ0X0MiXSxbNCw1LCIiLDAseyJzdHlsZSI6eyJ0YWlsIjp7Im5hbWUiOiJtYXBzIHRvIn19fV0sWzQsNiwiIiwyLHsic3R5bGUiOnsidGFpbCI6eyJuYW1lIjoibWFwcyB0byJ9fX1dLFs1LDcsIiIsMCx7InN0eWxlIjp7InRhaWwiOnsibmFtZSI6Im1hcHMgdG8ifX19XSxbNiw3LCIiLDIseyJzdHlsZSI6eyJ0YWlsIjp7Im5hbWUiOiJtYXBzIHRvIn19fV1d
\[\begin{tikzcd}
	{\mathcal{C}(D,A)} &&& {\mathcal{C}(C,A)} \\
	& g & {g\circ f} \\
	& {H\ g\ u} & {Hf(H\ g\ u)} \\
	{H\ D} &&& {H\ C}
	\arrow["{\mathcal{C}(f,A)}", from=1-1, to=1-4]
	\arrow["{t_D}"', from=1-1, to=4-1]
	\arrow["{H\ f}"', from=4-1, to=4-4]
	\arrow["{t_C}", from=1-4, to=4-4]
	\arrow[maps to, from=2-2, to=2-3]
	\arrow[maps to, from=2-2, to=3-2]
	\arrow[maps to, from=2-3, to=3-3]
	\arrow[maps to, from=3-2, to=3-3]
\end{tikzcd}\]
\[Hf(H\ g\ u) = (Hf\circ Hg)u = H(g\circ f) u\]
since $H$ is contravariant.

\[\phi(t) = t_A(1_A)\in H\ A\]
\[\psi(u) = t,\quad t_C(g) = H\ g\ u\]
\[\phi(\psi(u)) = \phi(t) = t_A(1_A) = H\ 1_A\ u = 1_{H\ A}\ u = u\]

\[\psi(\phi(s)) = \psi(s_A(1_A)) = t, \quad t_C(g) = (H\ g) (s_A(1_A))\]
\begin{equation}
\begin{aligned}
    t_C(g) &= (H\ g \circ s_A)\ 1_A \\
    &= s_C\ \mathcal{C}(g,A)\ 1_A\\
    &= s_C(g)
\end{aligned}
\end{equation}
For $g:C\rightarrow A$,
% https://q.uiver.app/#q=WzAsNCxbMCwwLCJcXG1hdGhjYWx7Q30oQSxBKSJdLFsxLDAsIlxcbWF0aGNhbHtDfShDLEEpIl0sWzAsMSwiSFxcIEEiXSxbMSwxLCJIXFwgQyJdLFswLDEsIlxcbWF0aGNhbHtDfShnLEEpIl0sWzAsMiwic19BIiwyXSxbMiwzLCJIXFwgZyIsMl0sWzEsMywiY19BIl1d
\[\begin{tikzcd}
	{\mathcal{C}(A,A)} & {\mathcal{C}(C,A)} \\
	{H\ A} & {H\ C}
	\arrow["{\mathcal{C}(g,A)}", from=1-1, to=1-2]
	\arrow["{s_A}"', from=1-1, to=2-1]
	\arrow["{H\ g}"', from=2-1, to=2-2]
	\arrow["{c_A}", from=1-2, to=2-2]
\end{tikzcd}\]

\end{proof}

\begin{corollary}
    Yoneda's Lemma [\ref{yo}] $\implies$ Faux-neda [\ref{faux}].
\end{corollary}
\begin{example}
    \[\phi : \text{Nat}(\mathcal{C}(-,A), H)\cong H\ A\]
    If $H = \mathcal{C}(-,B)$, then $\psi(f) : \mathcal{C}(-,A)\rightarrow \mathcal{C}(-,B)$,
    \[\psi(f)_C (g) = \mathcal{C}(g,B)(f) = f\circ g\]
    so $\psi(f) = \mathcal{C}(A,f)$
    \[\phi(t) = \mathcal{C}(A,f)(1_A) = f\]
\end{example}

\begin{fact}
    Yoneda embedding preserves all limits.
\end{fact}

\end{document}
